\documentclass[11pt,a4paper]{article}
\usepackage[utf8]{inputenc}
\usepackage[T1]{fontenc}
\usepackage[spanish,es-minimal,es-noshorthands]{babel}
\AtBeginDocument{\renewcommand{\sin}{\operatorname{sen}}\renewcommand{\arcsin}{\operatorname{arcsen}}}
\usepackage{lmodern}
\usepackage[sfdefault,lining]{FiraSans}
\usepackage[a4paper,top=2.3cm,bottom=2.2cm,left=1.8cm,right=1.8cm,headheight=1.3cm,headsep=0.4cm,footskip=1cm]{geometry}
\usepackage{amsmath,amssymb}
\usepackage{icomma}
\usepackage{graphicx}
\usepackage[table]{xcolor}
\usepackage{tikz}
\usepackage{fontawesome5}
\usepackage{ccicons}
\usepackage{fancyhdr}
\usepackage{lastpage}
\usepackage{tasks}
\usepackage{needspace}
\usepackage[most]{tcolorbox}
\usepackage{array}
\usepackage[hidelinks]{hyperref}
\AtBeginDocument{\hypersetup{pdftitle={Ficha de ejercicios de la unidad 8: Áreas y perímetros. 1º ESO},pdfsubject={Matemáticas 1º ESO},pdfkeywords={perímetro, área, unidades de longitud, unidades de superficie, hectárea, cuadrado, rectángulo, romboide, triángulo, rombo, trapecio, polígonos regulares, apotema, circunferencia, número pi, círculo, sector circular, corona circular, figuras compuestas, teorema de Pitágoras, catetos, hipotenusa}}}

\definecolor{navy}{HTML}{1B2748}
\definecolor{accent}{HTML}{E8603A}
\definecolor{teal}{HTML}{1F9E90}
\definecolor{soft}{HTML}{EEF1F8}
\definecolor{muted}{HTML}{5D6685}

\setlength{\parindent}{0pt}
\setlength{\parskip}{3pt}
\settasks{label=\alph*), label-format=\bfseries\color{navy}, label-width=1.3em, item-indent=1.8em, column-sep=0.6em, after-item-skip=5pt, before-skip=3pt, after-skip=4pt}

% ---------- capçalera i peu ----------
\newcommand{\llicencia}{\href{https://creativecommons.org/licenses/by-nc-sa/4.0/deed.es}{\ccbyncsa}\ \ Licencia Creative Commons Atribución-NoComercial-CompartirIgual 4.0 Internacional (CC BY-NC-SA 4.0)}
\pagestyle{fancy}
\fancyhf{}
\fancyhead[L]{\raisebox{-0.35\height}{\includegraphics[height=0.95cm]{logo-gva.png}}}
\fancyhead[C]{\small\color{muted}Ficha de ejercicios · Unidad 8 · 1\textsuperscript{o} ESO}
\fancyhead[R]{\raisebox{-0.35\height}{\includegraphics[height=0.95cm]{logo-vermellar.png}}}
\fancyfoot[L]{\scriptsize\color{muted}\llicencia}
\fancyfoot[R]{\small\color{muted}\thepage}
\renewcommand{\headrulewidth}{0.4pt}
\renewcommand{\headrule}{\color{muted!40}\hrule height \headrulewidth}
\fancypagestyle{primera}{\fancyhf{}\renewcommand{\headrulewidth}{0pt}
  \fancyfoot[L]{\scriptsize\color{muted}\llicencia}
  \fancyfoot[R]{\small\color{muted}\thepage}}

% ---------- nivells i exercicis ----------
\newcommand{\estrelles}[1]{{\color{accent}\ifcase#1\or\faStar\faStar[regular]\faStar[regular]\or\faStar\faStar\faStar[regular]\or\faStar\faStar\faStar\fi}}
\newcounter{exer}
\newcommand{\ex}[1]{\par\needspace{4\baselineskip}\medskip\stepcounter{exer}\noindent{\color{navy}\bfseries\theexer.}\ \ {\small\estrelles{#1}}\ \ }

\newcommand{\seccio}[2]{\par\needspace{10\baselineskip}\bigskip
  \begin{tcolorbox}[colback=navy,colframe=navy,coltext=white,arc=2mm,boxrule=0pt,left=3mm,right=3mm,top=1.2mm,bottom=1.2mm]
  {\large\bfseries #1\quad #2}
  \end{tcolorbox}\vspace{-2pt}}

\newcolumntype{C}[1]{>{\centering\arraybackslash}p{#1}}

% quadre de recordatori (només a la fitxa d'exercicis)
\newcommand{\recorda}[1]{\par\smallskip\begin{tcolorbox}[colback=accent!7,colframe=accent!55,arc=2mm,boxrule=0.5pt,left=2.5mm,right=2.5mm,top=1mm,bottom=1mm]\small{\color{accent}\faLightbulb[regular]}\ \textbf{Recuerda:}\ #1\end{tcolorbox}}
\newcommand{\res}[1]{#1}\newcommand{\rest}[1]{#1}

\def\fmneed{14}
% =====================================================================
% Capçalera FentMates de les fitxes: un color per a cada curs.
% Els colors són els de la pàgina inicial (LEVELS a assets/js/materials.js);
% el logo ve de _eines/logo/pdf/ (python3 _eines/logo/genera.py).
%
% Ús, al final del preàmbul de la unitat:
%   % =====================================================================
% Capçalera FentMates de les fitxes: un color per a cada curs.
% Els colors són els de la pàgina inicial (LEVELS a assets/js/materials.js);
% el logo ve de _eines/logo/pdf/ (python3 _eines/logo/genera.py).
%
% Ús, al final del preàmbul de la unitat:
%   % =====================================================================
% Capçalera FentMates de les fitxes: un color per a cada curs.
% Els colors són els de la pàgina inicial (LEVELS a assets/js/materials.js);
% el logo ve de _eines/logo/pdf/ (python3 _eines/logo/genera.py).
%
% Ús, al final del preàmbul de la unitat:
%   \input{../comu/fentmates.tex}
%   \fmcurs{1eso}
% i, al començament del document:
%   \capcalera{Matemàtiques}{1r ESO}{Fitxa d'exercicis · Unitat 2}{Divisibilitat}
% Opcional abans de \input: \def\fmneed{14} (espai mínim abans de cada secció).
% =====================================================================
\makeatletter
\def\fm@logodir{../../logo/pdf}
\providecommand{\fmneed}{10}
\def\fm@cursnom{}

% clau del curs, color i color del text damunt del color
\newcommand{\fm@curs}[3]{\@namedef{fm@color@#1}{#2}\@namedef{fm@text@#1}{#3}}
\fm@curs{1eso}{E6B84A}{navy}
\fm@curs{2eso}{3B82F6}{white}
\fm@curs{3eso}{FF7A50}{white}
\fm@curs{4eso-a}{A78BFA}{white}
\fm@curs{4eso-b}{6D28D9}{white}
\fm@curs{1batx-mi}{9B1C1C}{white}
\fm@curs{1batx-ccss}{EF6B6B}{white}
\fm@curs{2batx-mii}{1E6B47}{white}
\fm@curs{2batx-ccss}{2BB3A3}{white}

\newcommand{\fmcurs}[1]{%
  \@ifundefined{fm@color@#1}{\PackageError{fentmates}{Curs desconegut: #1}{}}{}%
  \edef\fm@tmp{\noexpand\definecolor{curs}{HTML}{\@nameuse{fm@color@#1}}%
               \noexpand\colorlet{curstext}{\@nameuse{fm@text@#1}}}\fm@tmp
  \colorlet{cursfosc}{curs!72!black}%
  \def\fm@logo{\fm@logodir/vertical-#1.pdf}%
  \def\fm@marca{\fm@logodir/marca-#1.pdf}}

% ---------- capçalera de la primera pàgina ----------
% #1 matèria (petit), #2 curs, #3 tipus de document i unitat, #4 títol
\newcommand{\capcalera}[4]{%
  \gdef\fm@cursnom{#2}%
  \par\noindent
  \begin{tcolorbox}[colback=curs,colframe=curs,boxrule=0pt,arc=2.2mm,width=3.5cm,height=2.15cm,
      valign=center,left=3mm,right=2mm,top=1mm,bottom=1mm,nobeforeafter,box align=center]
    \color{curstext}{\footnotesize #1}\par\vspace{1pt}{\fontsize{19}{21}\selectfont\bfseries #2\par}%
  \end{tcolorbox}\hspace{5mm}%
  \begin{minipage}[c]{\dimexpr\linewidth-3.5cm-5mm-2.9cm\relax}
    {\color{cursfosc}\small\bfseries\MakeUppercase{#3}}\par\vspace{2pt}
    {\color{navy}\fontsize{20}{23}\selectfont\bfseries #4\par}
  \end{minipage}\hfill
  \begin{minipage}[c]{2.6cm}\raggedleft\includegraphics[width=2.6cm]{\fm@logo}\end{minipage}\par
  \vspace{2pt}}

% ---------- barres de secció del color del curs ----------
\renewcommand{\seccio}[2]{\par\needspace{\fmneed\baselineskip}\bigskip
  \begin{tcolorbox}[colback=curs,colframe=curs,coltext=curstext,arc=2mm,boxrule=0pt,left=3mm,right=3mm,top=1.2mm,bottom=1.2mm]
  {\large\bfseries #1\quad #2}
  \end{tcolorbox}\vspace{-2pt}}

% ---------- estrelles de dificultat del color del curs (a Batxillerat no se'n posen) ----------
\@ifundefined{estrelles}{}{%
  \renewcommand{\estrelles}[1]{{\color{curs}\ifcase#1\or\faStar\faStar[regular]\faStar[regular]\or\faStar\faStar\faStar[regular]\or\faStar\faStar\faStar\fi}}}

% ---------- capçalera i peu de les pàgines ----------
% primera pàgina: logos institucionals al peu
\fancypagestyle{primera}{\fancyhf{}\renewcommand{\headrulewidth}{0pt}
  \fancyfoot[L]{\raisebox{-0.32\height}{\includegraphics[height=0.78cm]{logo-gva.png}}\hspace{5mm}%
    \raisebox{-0.32\height}{\includegraphics[height=0.78cm]{logo-vermellar.png}}}
  \fancyfoot[R]{\scriptsize\color{muted}\href{https://creativecommons.org/licenses/by-nc-sa/4.0/deed.ca}{\ccbyncsa}\ \ CC BY-NC-SA 4.0 · fentmates.org\qquad\small\thepage}}
% la resta: el símbol a l'esquerra i el curs a la dreta, amb el títol al centre com abans
\fancyhead[L]{\raisebox{-0.3\height}{\includegraphics[height=0.62cm]{\fm@marca}}}
\fancyhead[R]{\small\bfseries\color{cursfosc}\fm@cursnom}
\renewcommand{\headrule}{\color{curs!45}\hrule height \headrulewidth}
\makeatother

%   \fmcurs{1eso}
% i, al començament del document:
%   \capcalera{Matemàtiques}{1r ESO}{Fitxa d'exercicis · Unitat 2}{Divisibilitat}
% Opcional abans de \input: \def\fmneed{14} (espai mínim abans de cada secció).
% =====================================================================
\makeatletter
\def\fm@logodir{../../logo/pdf}
\providecommand{\fmneed}{10}
\def\fm@cursnom{}

% clau del curs, color i color del text damunt del color
\newcommand{\fm@curs}[3]{\@namedef{fm@color@#1}{#2}\@namedef{fm@text@#1}{#3}}
\fm@curs{1eso}{E6B84A}{navy}
\fm@curs{2eso}{3B82F6}{white}
\fm@curs{3eso}{FF7A50}{white}
\fm@curs{4eso-a}{A78BFA}{white}
\fm@curs{4eso-b}{6D28D9}{white}
\fm@curs{1batx-mi}{9B1C1C}{white}
\fm@curs{1batx-ccss}{EF6B6B}{white}
\fm@curs{2batx-mii}{1E6B47}{white}
\fm@curs{2batx-ccss}{2BB3A3}{white}

\newcommand{\fmcurs}[1]{%
  \@ifundefined{fm@color@#1}{\PackageError{fentmates}{Curs desconegut: #1}{}}{}%
  \edef\fm@tmp{\noexpand\definecolor{curs}{HTML}{\@nameuse{fm@color@#1}}%
               \noexpand\colorlet{curstext}{\@nameuse{fm@text@#1}}}\fm@tmp
  \colorlet{cursfosc}{curs!72!black}%
  \def\fm@logo{\fm@logodir/vertical-#1.pdf}%
  \def\fm@marca{\fm@logodir/marca-#1.pdf}}

% ---------- capçalera de la primera pàgina ----------
% #1 matèria (petit), #2 curs, #3 tipus de document i unitat, #4 títol
\newcommand{\capcalera}[4]{%
  \gdef\fm@cursnom{#2}%
  \par\noindent
  \begin{tcolorbox}[colback=curs,colframe=curs,boxrule=0pt,arc=2.2mm,width=3.5cm,height=2.15cm,
      valign=center,left=3mm,right=2mm,top=1mm,bottom=1mm,nobeforeafter,box align=center]
    \color{curstext}{\footnotesize #1}\par\vspace{1pt}{\fontsize{19}{21}\selectfont\bfseries #2\par}%
  \end{tcolorbox}\hspace{5mm}%
  \begin{minipage}[c]{\dimexpr\linewidth-3.5cm-5mm-2.9cm\relax}
    {\color{cursfosc}\small\bfseries\MakeUppercase{#3}}\par\vspace{2pt}
    {\color{navy}\fontsize{20}{23}\selectfont\bfseries #4\par}
  \end{minipage}\hfill
  \begin{minipage}[c]{2.6cm}\raggedleft\includegraphics[width=2.6cm]{\fm@logo}\end{minipage}\par
  \vspace{2pt}}

% ---------- barres de secció del color del curs ----------
\renewcommand{\seccio}[2]{\par\needspace{\fmneed\baselineskip}\bigskip
  \begin{tcolorbox}[colback=curs,colframe=curs,coltext=curstext,arc=2mm,boxrule=0pt,left=3mm,right=3mm,top=1.2mm,bottom=1.2mm]
  {\large\bfseries #1\quad #2}
  \end{tcolorbox}\vspace{-2pt}}

% ---------- estrelles de dificultat del color del curs (a Batxillerat no se'n posen) ----------
\@ifundefined{estrelles}{}{%
  \renewcommand{\estrelles}[1]{{\color{curs}\ifcase#1\or\faStar\faStar[regular]\faStar[regular]\or\faStar\faStar\faStar[regular]\or\faStar\faStar\faStar\fi}}}

% ---------- capçalera i peu de les pàgines ----------
% primera pàgina: logos institucionals al peu
\fancypagestyle{primera}{\fancyhf{}\renewcommand{\headrulewidth}{0pt}
  \fancyfoot[L]{\raisebox{-0.32\height}{\includegraphics[height=0.78cm]{logo-gva.png}}\hspace{5mm}%
    \raisebox{-0.32\height}{\includegraphics[height=0.78cm]{logo-vermellar.png}}}
  \fancyfoot[R]{\scriptsize\color{muted}\href{https://creativecommons.org/licenses/by-nc-sa/4.0/deed.ca}{\ccbyncsa}\ \ CC BY-NC-SA 4.0 · fentmates.org\qquad\small\thepage}}
% la resta: el símbol a l'esquerra i el curs a la dreta, amb el títol al centre com abans
\fancyhead[L]{\raisebox{-0.3\height}{\includegraphics[height=0.62cm]{\fm@marca}}}
\fancyhead[R]{\small\bfseries\color{cursfosc}\fm@cursnom}
\renewcommand{\headrule}{\color{curs!45}\hrule height \headrulewidth}
\makeatother

%   \fmcurs{1eso}
% i, al començament del document:
%   \capcalera{Matemàtiques}{1r ESO}{Fitxa d'exercicis · Unitat 2}{Divisibilitat}
% Opcional abans de \input: \def\fmneed{14} (espai mínim abans de cada secció).
% =====================================================================
\makeatletter
\def\fm@logodir{../../logo/pdf}
\providecommand{\fmneed}{10}
\def\fm@cursnom{}

% clau del curs, color i color del text damunt del color
\newcommand{\fm@curs}[3]{\@namedef{fm@color@#1}{#2}\@namedef{fm@text@#1}{#3}}
\fm@curs{1eso}{E6B84A}{navy}
\fm@curs{2eso}{3B82F6}{white}
\fm@curs{3eso}{FF7A50}{white}
\fm@curs{4eso-a}{A78BFA}{white}
\fm@curs{4eso-b}{6D28D9}{white}
\fm@curs{1batx-mi}{9B1C1C}{white}
\fm@curs{1batx-ccss}{EF6B6B}{white}
\fm@curs{2batx-mii}{1E6B47}{white}
\fm@curs{2batx-ccss}{2BB3A3}{white}

\newcommand{\fmcurs}[1]{%
  \@ifundefined{fm@color@#1}{\PackageError{fentmates}{Curs desconegut: #1}{}}{}%
  \edef\fm@tmp{\noexpand\definecolor{curs}{HTML}{\@nameuse{fm@color@#1}}%
               \noexpand\colorlet{curstext}{\@nameuse{fm@text@#1}}}\fm@tmp
  \colorlet{cursfosc}{curs!72!black}%
  \def\fm@logo{\fm@logodir/vertical-#1.pdf}%
  \def\fm@marca{\fm@logodir/marca-#1.pdf}}

% ---------- capçalera de la primera pàgina ----------
% #1 matèria (petit), #2 curs, #3 tipus de document i unitat, #4 títol
\newcommand{\capcalera}[4]{%
  \gdef\fm@cursnom{#2}%
  \par\noindent
  \begin{tcolorbox}[colback=curs,colframe=curs,boxrule=0pt,arc=2.2mm,width=3.5cm,height=2.15cm,
      valign=center,left=3mm,right=2mm,top=1mm,bottom=1mm,nobeforeafter,box align=center]
    \color{curstext}{\footnotesize #1}\par\vspace{1pt}{\fontsize{19}{21}\selectfont\bfseries #2\par}%
  \end{tcolorbox}\hspace{5mm}%
  \begin{minipage}[c]{\dimexpr\linewidth-3.5cm-5mm-2.9cm\relax}
    {\color{cursfosc}\small\bfseries\MakeUppercase{#3}}\par\vspace{2pt}
    {\color{navy}\fontsize{20}{23}\selectfont\bfseries #4\par}
  \end{minipage}\hfill
  \begin{minipage}[c]{2.6cm}\raggedleft\includegraphics[width=2.6cm]{\fm@logo}\end{minipage}\par
  \vspace{2pt}}

% ---------- barres de secció del color del curs ----------
\renewcommand{\seccio}[2]{\par\needspace{\fmneed\baselineskip}\bigskip
  \begin{tcolorbox}[colback=curs,colframe=curs,coltext=curstext,arc=2mm,boxrule=0pt,left=3mm,right=3mm,top=1.2mm,bottom=1.2mm]
  {\large\bfseries #1\quad #2}
  \end{tcolorbox}\vspace{-2pt}}

% ---------- estrelles de dificultat del color del curs (a Batxillerat no se'n posen) ----------
\@ifundefined{estrelles}{}{%
  \renewcommand{\estrelles}[1]{{\color{curs}\ifcase#1\or\faStar\faStar[regular]\faStar[regular]\or\faStar\faStar\faStar[regular]\or\faStar\faStar\faStar\fi}}}

% ---------- capçalera i peu de les pàgines ----------
% primera pàgina: logos institucionals al peu
\fancypagestyle{primera}{\fancyhf{}\renewcommand{\headrulewidth}{0pt}
  \fancyfoot[L]{\raisebox{-0.32\height}{\includegraphics[height=0.78cm]{logo-gva.png}}\hspace{5mm}%
    \raisebox{-0.32\height}{\includegraphics[height=0.78cm]{logo-vermellar.png}}}
  \fancyfoot[R]{\scriptsize\color{muted}\href{https://creativecommons.org/licenses/by-nc-sa/4.0/deed.ca}{\ccbyncsa}\ \ CC BY-NC-SA 4.0 · fentmates.org\qquad\small\thepage}}
% la resta: el símbol a l'esquerra i el curs a la dreta, amb el títol al centre com abans
\fancyhead[L]{\raisebox{-0.3\height}{\includegraphics[height=0.62cm]{\fm@marca}}}
\fancyhead[R]{\small\bfseries\color{cursfosc}\fm@cursnom}
\renewcommand{\headrule}{\color{curs!45}\hrule height \headrulewidth}
\makeatother

\fmcurs{1eso}

\begin{document}
\thispagestyle{primera}

\capcalera{Matemáticas}{1\textsuperscript{o} ESO}{Ficha de ejercicios · Unidad 8}{Áreas y perímetros}
\vspace{10pt}
Nombre y apellidos: \rule{6.6cm}{0.4pt}\hfill Grupo: \rule{1.6cm}{0.4pt}\hfill Fecha: \rule{2.2cm}{0.4pt}
\vspace{8pt}

\begin{tcolorbox}[colback=soft,colframe=soft,arc=2mm,left=3mm,right=3mm,top=1.5mm,bottom=1.5mm]
\small Ejercicios graduados por niveles de dificultad:\quad \estrelles{1}\ básico \qquad \estrelles{2}\ medio \qquad \estrelles{3}\ avanzado.\\[2pt] Resuélvelos en el cuaderno mostrando todos los pasos y escribe siempre las unidades (cm, cm$^2$…). Cuando haya círculos o circunferencias, usa $\pi\approx3,14$. Si una raíz no es exacta, redondéala a las centésimas.
\end{tcolorbox}

% =====================================================
\seccio{1}{Perímetro y área. Unidades de medida}

\ex{1} Cada cuadradito mide 1\,cm de lado. Calcula el área y el perímetro de cada figura:\begin{tasks}(3)
  \task \begin{tikzpicture}[x=0.55cm,y=0.55cm,baseline=(current bounding box.center),line join=round,line cap=round]
\fill[teal!25] (0,0) rectangle (1,1);
\fill[teal!25] (1,0) rectangle (2,1);
\fill[teal!25] (2,0) rectangle (3,1);
\fill[teal!25] (0,1) rectangle (1,2);
\fill[teal!25] (1,1) rectangle (2,2);
\fill[teal!25] (2,1) rectangle (3,2);
\draw[muted!80,line width=0.5pt,step=1] (0.000,0.000) grid (3.000,2.000);
\draw[thick] (0.000,0.000) -- (1.000,0.000);
\draw[thick] (0.000,0.000) -- (0.000,1.000);
\draw[thick] (1.000,0.000) -- (2.000,0.000);
\draw[thick] (2.000,0.000) -- (3.000,0.000);
\draw[thick] (3.000,0.000) -- (3.000,1.000);
\draw[thick] (0.000,2.000) -- (1.000,2.000);
\draw[thick] (0.000,1.000) -- (0.000,2.000);
\draw[thick] (1.000,2.000) -- (2.000,2.000);
\draw[thick] (2.000,2.000) -- (3.000,2.000);
\draw[thick] (3.000,1.000) -- (3.000,2.000);
\end{tikzpicture}
  \task \begin{tikzpicture}[x=0.55cm,y=0.55cm,baseline=(current bounding box.center),line join=round,line cap=round]
\fill[teal!25] (0,0) rectangle (1,1);
\fill[teal!25] (1,0) rectangle (2,1);
\fill[teal!25] (2,0) rectangle (3,1);
\fill[teal!25] (3,0) rectangle (4,1);
\fill[teal!25] (3,1) rectangle (4,2);
\fill[teal!25] (3,2) rectangle (4,3);
\fill[teal!25] (0,1) rectangle (1,2);
\draw[muted!80,line width=0.5pt,step=1] (0.000,0.000) grid (4.000,3.000);
\draw[thick] (0.000,0.000) -- (1.000,0.000);
\draw[thick] (0.000,0.000) -- (0.000,1.000);
\draw[thick] (1.000,0.000) -- (2.000,0.000);
\draw[thick] (1.000,1.000) -- (2.000,1.000);
\draw[thick] (2.000,0.000) -- (3.000,0.000);
\draw[thick] (2.000,1.000) -- (3.000,1.000);
\draw[thick] (3.000,0.000) -- (4.000,0.000);
\draw[thick] (4.000,0.000) -- (4.000,1.000);
\draw[thick] (3.000,1.000) -- (3.000,2.000);
\draw[thick] (4.000,1.000) -- (4.000,2.000);
\draw[thick] (3.000,3.000) -- (4.000,3.000);
\draw[thick] (3.000,2.000) -- (3.000,3.000);
\draw[thick] (4.000,2.000) -- (4.000,3.000);
\draw[thick] (0.000,2.000) -- (1.000,2.000);
\draw[thick] (0.000,1.000) -- (0.000,2.000);
\draw[thick] (1.000,1.000) -- (1.000,2.000);
\end{tikzpicture}
  \task \begin{tikzpicture}[x=0.55cm,y=0.55cm,baseline=(current bounding box.center),line join=round,line cap=round]
\fill[teal!25] (1,0) rectangle (2,1);
\fill[teal!25] (0,1) rectangle (1,2);
\fill[teal!25] (1,1) rectangle (2,2);
\fill[teal!25] (2,1) rectangle (3,2);
\fill[teal!25] (1,2) rectangle (2,3);
\draw[muted!80,line width=0.5pt,step=1] (0.000,0.000) grid (3.000,3.000);
\draw[thick] (1.000,0.000) -- (2.000,0.000);
\draw[thick] (1.000,0.000) -- (1.000,1.000);
\draw[thick] (2.000,0.000) -- (2.000,1.000);
\draw[thick] (0.000,1.000) -- (1.000,1.000);
\draw[thick] (0.000,2.000) -- (1.000,2.000);
\draw[thick] (0.000,1.000) -- (0.000,2.000);
\draw[thick] (2.000,1.000) -- (3.000,1.000);
\draw[thick] (2.000,2.000) -- (3.000,2.000);
\draw[thick] (3.000,1.000) -- (3.000,2.000);
\draw[thick] (1.000,3.000) -- (2.000,3.000);
\draw[thick] (1.000,2.000) -- (1.000,3.000);
\draw[thick] (2.000,2.000) -- (2.000,3.000);
\end{tikzpicture}
\end{tasks}\recorda{El \textbf{perímetro} es la longitud del contorno (cm, m…) y el \textbf{área} es la superficie que ocupa la figura (cm$^2$, m$^2$…). En una cuadrícula, el área es el número de cuadraditos y el perímetro, el número de lados de cuadradito del contorno.}

\ex{1} Expresa en metros:\begin{tasks}(5)
  \task 4,7\,km
  \task 380\,cm
  \task 25\,dm
  \task 9\,hm
  \task 1250\,mm
\end{tasks}\recorda{Las unidades de longitud van de 10 en 10: km, hm, dam, m, dm, cm, mm. Para pasar a una unidad más pequeña multiplicamos y para pasar a una más grande dividimos.}

\ex{1} Haz los cambios de unidades:\begin{tasks}(4)
  \task 3,2\,m a cm
  \task 45\,mm a cm
  \task 0,8\,km a dam
  \task 760\,dm a hm
\end{tasks}

\ex{1} Haz los cambios de unidades de superficie:\begin{tasks}(3)
  \task 5\,dm$^2$ a cm$^2$
  \task 2,4\,m$^2$ a dm$^2$
  \task 35\,000\,cm$^2$ a m$^2$
  \task 0,7\,hm$^2$ a m$^2$
  \task 8\,km$^2$ a hm$^2$
\end{tasks}\recorda{Las unidades de superficie van de 100 en 100: \ $1\text{ m}^2=100\text{ dm}^2=10\,000\text{ cm}^2$.}

\ex{2} Haz los cambios con unidades agrarias ($1\text{ ha}=1\text{ hm}^2$, \ $1\text{ a}=1\text{ dam}^2$):\begin{tasks}(4)
  \task 3,5\,ha a m$^2$
  \task 4500\,m$^2$ a ha
  \task 12\,a a m$^2$
  \task 0,25\,km$^2$ a ha
\end{tasks}

\ex{2} Ordena de menor a mayor: \ $0,05\text{ km}^2$, \ $4\text{ ha}$, \ $520\text{ dam}^2$, \ $48\,000\text{ m}^2$.

\ex{2} Una finca de 2,4\,ha se vende a 18\,€ el metro cuadrado. ¿Cuánto cuesta?

\ex{3} Queremos hacer rectángulos de 20\,cm de perímetro con lados que sean números naturales. Escríbelos todos y calcula su área. ¿Cuál tiene el área más grande?


% =====================================================
\seccio{2}{Perímetro de los polígonos}

\ex{1} Calcula el perímetro de:\begin{tasks}(2)
  \task un triángulo de lados 7,5\,cm, 9\,cm y 11\,cm
  \task un cuadrado de 8,5\,cm de lado
  \task un rectángulo de 15\,m × 9\,m
  \task un rombo de 6,3\,cm de lado
\end{tasks}\recorda{El perímetro de un polígono es la suma de sus lados. Cuadrado: $P=4\cdot l$. \ Rectángulo: $P=2\cdot b+2\cdot h$. \ Polígono regular de $n$ lados: $P=n\cdot l$.}

\ex{1} Calcula el perímetro de estos polígonos regulares:\begin{tasks}(2)
  \task pentágono de 7\,cm de lado
  \task hexágono de 4,5\,cm de lado
  \task octógono de 3,2\,cm de lado
  \task decágono de 12\,mm de lado
\end{tasks}

\ex{1} Calcula el perímetro de este polígono:\par\smallskip\begin{center}\begin{tikzpicture}[x=0.45cm,y=0.45cm,baseline=(current bounding box.center),line join=round,line cap=round]
\draw[fill=teal!12,draw=teal,thick] (0.000,0.000) -- (5.000,0.000) -- (5.000,3.000) -- (2.000,5.500) -- (0.000,3.000) -- cycle;
\node[font=\small,below] at (2.500,0.000) {5 cm};
\node[font=\small,right] at (5.000,1.500) {3 cm};
\node[font=\small,above right] at (3.500,4.250) {4 cm};
\node[font=\small,above left] at (1.000,4.250) {3,5 cm};
\node[font=\small,left] at (0.000,1.500) {3 cm};
\end{tikzpicture}\end{center}

\ex{2} Calcula el lado de:\begin{tasks}(1)
  \task un cuadrado de 46\,cm de perímetro
  \task un hexágono regular de 51\,cm de perímetro
  \task un triángulo equilátero de 13,5\,m de perímetro
\end{tasks}

\ex{2} Un rectángulo tiene 64\,cm de perímetro y la base mide 20\,cm. ¿Cuánto mide la altura?

\ex{2} Queremos cercar un campo rectangular de 120\,m × 85\,m con tres hilos de alambre. ¿Cuántos metros de alambre hacen falta? ¿Cuánto costarán si el metro vale 0,35\,€?

\ex{3} La base de un rectángulo mide el triple que la altura y el perímetro mide 96\,cm. ¿Cuánto miden la base y la altura?

\ex{3} Cada uno de los dos lados iguales de un triángulo isósceles mide el doble que la base. Si el perímetro es de 40\,cm, ¿cuánto mide cada lado?


% =====================================================
\seccio{3}{Área del cuadrado y del rectángulo}

\ex{1} Calcula el área y el perímetro de:\begin{tasks}(2)
  \task un cuadrado de 12\,cm de lado
  \task un rectángulo de 15\,cm × 8\,cm
  \task un rectángulo de 7,5\,m × 4\,m
  \task un cuadrado de 2,5\,m de lado
\end{tasks}\recorda{Rectángulo: $A=b\cdot h$. \ Cuadrado: $A=l\cdot l=l^2$. Las dos medidas tienen que estar en la misma unidad.}

\ex{1} Un rectángulo tiene 84\,cm$^2$ de área y 12\,cm de base. ¿Cuánto miden la altura y el perímetro?

\ex{1} Responde:\begin{tasks}(1)
  \task Un cuadrado tiene 121\,m$^2$ de área. ¿Cuánto mide el lado? ¿Y el perímetro?
  \task Un cuadrado tiene 28\,cm de perímetro. ¿Cuál es su área?
\end{tasks}

\ex{2} Una habitación mide 5,4\,m × 3,5\,m. ¿Cuánto costará poner parqué si el metro cuadrado vale 24\,€?

\ex{2} Queremos pintar una pared de 6\,m × 2,5\,m que tiene una ventana de 1,5\,m × 1,2\,m. Un bote de pintura cubre 5\,m$^2$. ¿Cuántos botes tenemos que comprar?

\ex{2} ¿Cuántas baldosas cuadradas de 50\,cm de lado hacen falta para embaldosar un patio de 8\,m × 6\,m?

\ex{3} Un cuadrado mide 5\,cm de lado. Si duplicamos el lado, ¿qué le pasa al perímetro? ¿Y al área?

\ex{3} Escribe todos los rectángulos de lados naturales que tienen 36\,cm$^2$ de área. ¿Cuál tiene el perímetro más pequeño?


% =====================================================
\seccio{4}{Área del romboide y del triángulo}

\ex{1} Calcula el área de los romboides:\begin{tasks}(2)
  \task base 12\,cm y altura 7\,cm
  \task base 8,5\,m y altura 4\,m
\end{tasks}\recorda{Romboide: $A=b\cdot h$. \ Triángulo: $A=\dfrac{b\cdot h}{2}$. La altura siempre es perpendicular a la base.}

\ex{1} Calcula el área de los triángulos:\begin{tasks}(3)
  \task base 14\,cm y altura 9\,cm
  \task base 7\,cm y altura 5\,cm
  \task base 6,4\,m y altura 3,5\,m
\end{tasks}

\ex{1} Calcula el área de cada triángulo:\begin{tasks}(3)
  \task \begin{tikzpicture}[x=0.42cm,y=0.42cm,baseline=(current bounding box.center),line join=round,line cap=round]
\draw[fill=teal!12,draw=teal,thick] (0.000,0.000) -- (6.000,0.000) -- (2.000,4.000) -- cycle;
\draw[dashed,muted] (2.000,0.000) -- (2.000,4.000);
\draw[thin] (2.524,0.000) -- (2.524,0.524) -- (2.000,0.524);
\node[font=\small,below] at (3.000,0.000) {6 cm};
\node[font=\small,right] at (2.000,2.000) {4 cm};
\end{tikzpicture}
  \task \begin{tikzpicture}[x=0.38cm,y=0.38cm,baseline=(current bounding box.center),line join=round,line cap=round]
\draw[fill=teal!12,draw=teal,thick] (0.000,0.000) -- (8.000,0.000) -- (0.000,5.000) -- cycle;
\draw[thin] (0.579,0.000) -- (0.579,0.579) -- (0.000,0.579);
\node[font=\small,below] at (4.000,0.000) {8 cm};
\node[font=\small,left] at (0.000,2.500) {5 cm};
\end{tikzpicture}
  \task \begin{tikzpicture}[x=0.38cm,y=0.38cm,baseline=(current bounding box.center),line join=round,line cap=round]
\draw[fill=teal!12,draw=teal,thick] (0.000,0.000) -- (4.000,0.000) -- (6.500,5.000) -- cycle;
\draw[dashed,muted] (4.000,0.000) -- (6.500,0.000);
\draw[dashed,muted] (6.500,0.000) -- (6.500,5.000);
\draw[thin] (5.921,0.000) -- (5.921,0.579) -- (6.500,0.579);
\node[font=\small,below] at (2.000,0.000) {4 cm};
\node[font=\small,right] at (6.500,2.500) {5 cm};
\end{tikzpicture}
\end{tasks}

\ex{2} Los catetos de un triángulo rectángulo miden 9\,cm y 12\,cm, y la hipotenusa, 15\,cm. Calcula su área y su perímetro.

\ex{2} Responde:\begin{tasks}(1)
  \task Un romboide tiene 96\,cm$^2$ de área y 12\,cm de base. ¿Cuánto mide la altura?
  \task Un triángulo tiene 45\,m$^2$ de área y 9\,m de altura. ¿Cuánto mide la base?
\end{tasks}

\ex{2} Cada cuadradito de la cuadrícula mide 1\,cm de lado. Calcula el área del triángulo.\par\smallskip\begin{center}\begin{tikzpicture}[x=0.42cm,y=0.42cm,baseline=(current bounding box.center),line join=round,line cap=round]
\draw[muted!80,line width=0.5pt,step=1] (0.000,0.000) grid (8.000,5.000);
\draw[fill=accent,fill opacity=0.22,draw=accent,thick] (0.000,0.000) -- (8.000,0.000) -- (3.000,5.000) -- cycle;
\end{tikzpicture}\end{center}

\ex{3} Un triángulo tiene los lados de 13\,cm, 14\,cm y 15\,cm, y la altura sobre el lado de 14\,cm mide 12\,cm. Calcula su área y la altura sobre el lado de 15\,cm.

\ex{3} Dibuja dos rectas paralelas separadas 6\,cm y, sobre una de ellas, un segmento de 10\,cm. Dibuja tres triángulos diferentes que tengan este segmento como base y el tercer vértice sobre la otra recta. ¿Qué área tiene cada uno? ¿Por qué?


% =====================================================
\seccio{5}{Área del rombo y del trapecio}

\ex{1} Calcula el área de los rombos:\begin{tasks}(3)
  \task diagonales 16\,cm y 10\,cm
  \task diagonales 9\,m y 6\,m
  \task diagonales 7,5\,cm y 4\,cm
\end{tasks}\recorda{Rombo: $A=\dfrac{D\cdot d}{2}$. \ Trapecio: $A=\dfrac{(B+b)\cdot h}{2}$.}

\ex{1} Calcula el área de los trapecios:\begin{tasks}(2)
  \task bases 12\,cm y 8\,cm, altura 5\,cm
  \task bases 9\,m y 4\,m, altura 6\,m
  \task bases 10,5\,cm y 6,5\,cm, altura 4\,cm
\end{tasks}

\ex{2} Un rombo tiene 52\,cm de perímetro y las diagonales miden 24\,cm y 10\,cm. ¿Cuánto mide el lado? Calcula su área.

\ex{2} Calcula el área y el perímetro de este trapecio isósceles:\par\smallskip\begin{center}\begin{tikzpicture}[x=0.3cm,y=0.3cm,baseline=(current bounding box.center),line join=round,line cap=round]
\draw[fill=teal!12,draw=teal,thick] (0.000,0.000) -- (18.000,0.000) -- (14.000,3.000) -- (4.000,3.000) -- cycle;
\draw[dashed,muted] (4.000,0.000) -- (4.000,3.000);
\draw[thin] (4.733,0.000) -- (4.733,0.733) -- (4.000,0.733);
\node[font=\small,below] at (9.000,0.000) {18 cm};
\node[font=\small,above] at (9.000,3.000) {10 cm};
\node[font=\small,right] at (4.000,1.500) {3 cm};
\node[font=\small,right] at (16.000,1.500) {5 cm};
\end{tikzpicture}\end{center}

\ex{2} Responde:\begin{tasks}(1)
  \task Un rombo tiene 60\,cm$^2$ de área y la diagonal mayor mide 15\,cm. ¿Cuánto mide la menor?
  \task Un trapecio tiene 72\,m$^2$ de área y las bases miden 14\,m y 10\,m. ¿Cuánto mide la altura?
\end{tasks}

\ex{3} Una cometa tiene forma de rombo, con diagonales de 1,2\,m y 80\,cm. ¿Cuántos m$^2$ de tela hacen falta para hacerla? Exprésalo también en cm$^2$.

\ex{3} Una parcela tiene forma de trapecio, con bases de 85\,m y 55\,m y 40\,m de anchura (altura). ¿Cuántas hectáreas mide? ¿Cuánto vale si el metro cuadrado se paga a 12\,€?


% =====================================================
\seccio{6}{Área de los polígonos regulares y de polígonos cualesquiera}

\ex{1} Un pentágono regular tiene 10\,cm de lado y 6,88\,cm de apotema. Calcula su perímetro y su área.\recorda{Polígono regular: $A=\dfrac{P\cdot a}{2}$, donde $P=n\cdot l$ es el perímetro y $a$ es la apotema (el segmento del centro al punto medio de un lado).}

\ex{1} Un hexágono regular tiene 8\,cm de lado y 6,93\,cm de apotema. Calcula su perímetro y su área.

\ex{2} Un octógono regular tiene 40\,cm de perímetro y 6,04\,cm de apotema. ¿Cuánto mide el lado? ¿Y el área?

\ex{2} Un hexágono regular tiene 36\,cm de perímetro y 93,6\,cm$^2$ de área. ¿Cuánto mide la apotema?

\ex{2} Cada cuadradito mide 1\,cm de lado. Calcula el área del polígono descomponiéndolo en un rectángulo y un triángulo.\par\smallskip\begin{center}\begin{tikzpicture}[x=0.4cm,y=0.4cm,baseline=(current bounding box.center),line join=round,line cap=round]
\draw[muted!80,line width=0.5pt,step=1] (0.000,0.000) grid (6.000,6.000);
\draw[fill=accent,fill opacity=0.22,draw=accent,thick] (0.000,0.000) -- (6.000,0.000) -- (6.000,4.000) -- (3.000,6.000) -- (0.000,4.000) -- cycle;
\end{tikzpicture}\end{center}

\ex{3} Cada cuadradito mide 1\,cm de lado. Calcula el área del triángulo completándolo hasta un rectángulo y restando los triángulos que sobran.\par\smallskip\begin{center}\begin{tikzpicture}[x=0.4cm,y=0.4cm,baseline=(current bounding box.center),line join=round,line cap=round]
\draw[muted!80,line width=0.5pt,step=1] (0.000,0.000) grid (7.000,6.000);
\draw[fill=accent,fill opacity=0.22,draw=accent,thick] (0.000,0.000) -- (7.000,2.000) -- (3.000,6.000) -- cycle;
\end{tikzpicture}\end{center}

\ex{3} Queremos embaldosar 25\,m$^2$ con baldosas en forma de hexágono regular de 10\,cm de lado y 8,66\,cm de apotema. ¿Cuántas baldosas hacen falta, como mínimo?


% =====================================================
\seccio{7}{La circunferencia y su longitud}

\ex{1} ¿Cómo se llaman los elementos $a$, $b$, $c$ y $d$ de la figura?\par\smallskip\begin{center}\begin{tikzpicture}[x=0.75cm,y=0.75cm,baseline=(current bounding box.center),line join=round,line cap=round]
\draw[draw=teal,thick] (0.000,0.000) circle (2);
\draw[thick] (0.000,0.000) -- (-1.286,1.532);
\draw[thick] (-1.879,-0.684) -- (1.879,0.684);
\draw[thick] (-0.684,-1.879) -- (1.532,-1.286);
\draw[accent,very thick] (1.532,1.286) arc[start angle=40,end angle=100,radius=2];
\fill (0.000,0.000) circle (1.5pt);
\node[font=\small,inner sep=1pt] at (0.250,-0.300) {$O$};
\node[font=\small,above right] at (-0.643,0.766) {$a$};
\node[font=\small,below right] at (0.000,0.000) {$b$};
\node[font=\small,above] at (0.424,-1.582) {$c$};
\node[font=\small,inner sep=1pt] at (0.804,2.208) {$d$};
\end{tikzpicture}\end{center}

\ex{1} Calcula la longitud de las circunferencias (usa $\pi\approx3,14$):\begin{tasks}(4)
  \task radio 4\,cm
  \task radio 9\,m
  \task diámetro 20\,cm
  \task diámetro 7\,mm
\end{tasks}\recorda{La longitud de la circunferencia es $L=\pi\cdot d=2\cdot\pi\cdot r$, con $\pi\approx3,14$. Un arco de $n^\circ$ mide $L_{\text{arco}}=\dfrac{2\cdot\pi\cdot r\cdot n}{360}$.}

\ex{2} Responde:\begin{tasks}(1)
  \task Una circunferencia mide 31,4\,cm de longitud. ¿Cuánto mide el radio?
  \task Una circunferencia mide 75,36\,m de longitud. ¿Cuánto miden el diámetro y el radio?
\end{tasks}

\ex{2} La rueda de una bicicleta mide 70\,cm de diámetro. ¿Cuántos centímetros avanza en una vuelta? ¿Cuántas vueltas tiene que dar, aproximadamente, para recorrer 1\,km?

\ex{2} Calcula la longitud de los arcos:\begin{tasks}(3)
  \task $90^\circ$ y radio 8\,cm
  \task $60^\circ$ y radio 12\,cm
  \task $150^\circ$ y radio 6\,cm
\end{tasks}

\ex{3} Calcula el perímetro de:\begin{tasks}(2)
  \task un semicírculo de 10\,cm de diámetro
  \task un cuarto de círculo de 4\,cm de radio
\end{tasks}

\ex{3} La aguja de los minutos de un reloj mide 15\,cm. ¿Qué distancia recorre la punta en 20 minutos? ¿Y en una hora?


% =====================================================
\seccio{8}{El área del círculo y de las figuras circulares}

\ex{1} Calcula el área de los círculos ($\pi\approx3,14$):\begin{tasks}(4)
  \task radio 6\,cm
  \task radio 2,5\,m
  \task diámetro 14\,cm
  \task diámetro 30\,cm
\end{tasks}\recorda{Círculo: $A=\pi\cdot r^2$. \ Sector circular de $n^\circ$: $A=\dfrac{\pi\cdot r^2\cdot n}{360}$. \ Corona circular: $A=\pi\cdot R^2-\pi\cdot r^2$. Si te dan el diámetro, calcula primero el radio.}

\ex{2} Calcula el área de:\begin{tasks}(2)
  \task un semicírculo de 10\,cm de radio
  \task un cuarto de círculo de 6\,cm de radio
\end{tasks}

\ex{2} Calcula el área de los sectores circulares:\begin{tasks}(3)
  \task $45^\circ$ y radio 8\,cm
  \task $120^\circ$ y radio 9\,cm
  \task $270^\circ$ y radio 4\,cm
\end{tasks}

\ex{2} Calcula el área de la corona circular de la figura en cada caso:\par\smallskip\begin{center}\begin{tikzpicture}[x=0.6cm,y=0.6cm,baseline=(current bounding box.center),line join=round,line cap=round]
\draw[fill=accent!22,draw=accent,thick] (0.000,0.000) circle (2.1);
\draw[fill=white,draw=teal,thick] (0.000,0.000) circle (1.2);
\fill (0.000,0.000) circle (1.5pt);
\draw[thick] (0.000,0.000) -- (1.819,1.050);
\draw[dashed,muted] (0.000,0.000) -- (-1.128,0.410);
\node[font=\small,inner sep=1pt] at (1.226,1.104) {$R$};
\node[font=\small,inner sep=1pt] at (-0.600,0.021) {$r$};
\end{tikzpicture}\end{center}\begin{tasks}(2)
  \task $R=7$\,cm, \ $r=4$\,cm
  \task $R=10$\,m, \ $r=6$\,m
\end{tasks}

\ex{2} Responde:\begin{tasks}(1)
  \task Un círculo tiene 78,5\,cm$^2$ de área. ¿Cuánto mide el radio?
  \task Una circunferencia mide 18,84\,cm. ¿Cuál es el área del círculo que encierra?
\end{tasks}

\ex{3} Un CD mide 12\,cm de diámetro y tiene un agujero central de 1,5\,cm de diámetro. ¿Cuál es la superficie del CD?

\ex{3} ¿Qué tiene más superficie: una pizza de 32\,cm de diámetro o dos pizzas de 22\,cm de diámetro?


% =====================================================
\seccio{9}{Figuras compuestas}

\ex{1} Calcula el área y el perímetro de esta figura:\par\smallskip\begin{center}\begin{tikzpicture}[x=0.32cm,y=0.32cm,baseline=(current bounding box.center),line join=round,line cap=round]
\draw[fill=teal!12,draw=teal,thick] (0.000,0.000) -- (10.000,0.000) -- (10.000,4.000) -- (6.000,4.000) -- (6.000,7.000) -- (0.000,7.000) -- cycle;
\node[font=\small,below] at (5.000,0.000) {10 cm};
\node[font=\small,left] at (0.000,3.500) {7 cm};
\node[font=\small,right] at (10.000,2.000) {4 cm};
\node[font=\small,above] at (8.000,4.000) {4 cm};
\node[font=\small,above] at (3.000,7.000) {6 cm};
\end{tikzpicture}\end{center}\recorda{Para el área de una figura compuesta, la \textbf{descomponemos} en figuras conocidas y sumamos, o la \textbf{completamos} y restamos. El perímetro es solo el contorno exterior.}

\ex{1} Calcula el área de la fachada de esta casita:\par\smallskip\begin{center}\begin{tikzpicture}[x=0.3cm,y=0.3cm,baseline=(current bounding box.center),line join=round,line cap=round]
\draw[fill=teal!12,draw=teal,thick] (0.000,0.000) -- (8.000,0.000) -- (8.000,5.000) -- (0.000,5.000) -- cycle;
\draw[fill=accent!22,draw=accent,thick] (0.000,5.000) -- (8.000,5.000) -- (4.000,8.000) -- cycle;
\draw[dashed,muted] (4.000,5.000) -- (4.000,8.000);
\node[font=\small,below] at (4.000,0.000) {8 m};
\node[font=\small,right] at (8.000,2.500) {5 m};
\node[font=\small,right] at (4.000,6.500) {3 m};
\end{tikzpicture}\end{center}

\ex{2} Calcula el área y el perímetro de la figura, formada por un rectángulo y un semicírculo ($\pi\approx3,14$):\par\smallskip\begin{center}\begin{tikzpicture}[x=0.38cm,y=0.38cm,baseline=(current bounding box.center),line join=round,line cap=round]
\draw[fill=accent!22,draw=accent,thick] (0.000,0.000) -- (6.000,0.000) -- (6.000,4.000) -- (0.000,4.000) -- cycle;
\draw[fill=accent!22,draw=accent,thick] (6.000,2.000) -- (6.000,0.000) arc[start angle=-90,end angle=90,radius=2] -- cycle;
\draw[dashed,muted] (6.000,0.000) -- (6.000,4.000);
\node[font=\small,below] at (3.000,0.000) {6 cm};
\node[font=\small,left] at (0.000,2.000) {4 cm};
\end{tikzpicture}\end{center}

\ex{2} Calcula el área de la zona coloreada (un cuadrado de 8\,cm de lado sin cuatro cuartos de círculo de 4\,cm de radio):\par\smallskip\begin{center}\begin{tikzpicture}[x=0.28cm,y=0.28cm,baseline=(current bounding box.center),line join=round,line cap=round]
\draw[fill=accent!22,draw=accent,thick] (0.000,0.000) -- (8.000,0.000) -- (8.000,8.000) -- (0.000,8.000) -- cycle;
\draw[fill=white,draw=teal,thick] (0.000,0.000) -- (4.000,0.000) arc[start angle=0,end angle=90,radius=4] -- cycle;
\draw[fill=white,draw=teal,thick] (8.000,0.000) -- (8.000,4.000) arc[start angle=90,end angle=180,radius=4] -- cycle;
\draw[fill=white,draw=teal,thick] (8.000,8.000) -- (4.000,8.000) arc[start angle=180,end angle=270,radius=4] -- cycle;
\draw[fill=white,draw=teal,thick] (0.000,8.000) -- (-0.000,4.000) arc[start angle=270,end angle=360,radius=4] -- cycle;
\node[font=\small,below] at (4.000,0.000) {8 cm};
\end{tikzpicture}\end{center}

\ex{2} Alrededor de una piscina de 10\,m × 6\,m se hace una acera de 1,5\,m de anchura. ¿Cuál es el área de la acera?\par\smallskip\begin{center}\begin{tikzpicture}[x=0.3cm,y=0.3cm,baseline=(current bounding box.center),line join=round,line cap=round]
\draw[fill=accent!22,draw=accent,thick] (-1.500,-1.500) -- (11.500,-1.500) -- (11.500,7.500) -- (-1.500,7.500) -- cycle;
\draw[fill=white,draw=teal,thick] (0.000,0.000) -- (10.000,0.000) -- (10.000,6.000) -- (0.000,6.000) -- cycle;
\node[font=\small,inner sep=1pt] at (5.000,3.000) {10 m × 6 m};
\node[font=\small,inner sep=1pt] at (5.000,6.750) {1,5 m};
\end{tikzpicture}\end{center}

\ex{2} Una pista está formada por un rectángulo de 80\,m × 50\,m y dos semicírculos de 50\,m de diámetro en los extremos. Calcula su área y la longitud de una vuelta.

\ex{3} Una diana tiene tres circunferencias concéntricas de radios 3, 6 y 9\,cm. Calcula el área de cada zona.\par\smallskip\begin{center}\begin{tikzpicture}[x=0.55cm,y=0.55cm,baseline=(current bounding box.center),line join=round,line cap=round]
\draw[fill=accent!22,draw=accent,thick] (0.000,0.000) circle (2.7);
\draw[fill=white,draw=teal,thick] (0.000,0.000) circle (1.8);
\draw[fill=accent!8,draw=accent!60] (0.000,0.000) circle (1.8);
\draw[fill=accent!22,draw=accent,thick] (0.000,0.000) circle (0.9);
\fill (0.000,0.000) circle (1.5pt);
\draw[thick] (0.000,0.000) -- (2.700,0.000);
\node[font=\small,above] at (0.450,0.000) {3};
\node[font=\small,above] at (1.350,0.000) {3};
\node[font=\small,above] at (2.250,0.000) {3};
\end{tikzpicture}\end{center}

\ex{3} Dentro de un cuadrado de 6\,cm de lado se dibujan dos cuartos de circunferencia con centro en dos vértices opuestos. Calcula el área de la hoja que queda entre los dos arcos.\par\smallskip\begin{center}\begin{tikzpicture}[x=0.38cm,y=0.38cm,baseline=(current bounding box.center),line join=round,line cap=round]
\draw[thick,muted] (0.000,0.000) -- (6.000,0.000) -- (6.000,6.000) -- (0.000,6.000) -- cycle;
\draw[accent,very thick] (6.000,0.000) arc[start angle=0,end angle=90,radius=6];
\draw[accent,very thick] (0.000,6.000) arc[start angle=180,end angle=270,radius=6];
\draw[dashed,muted] (6.000,0.000) -- (0.000,6.000);
\node[font=\small,below] at (3.000,0.000) {6 cm};
\end{tikzpicture}\end{center}

\ex{3} Una ventana está formada por un rectángulo de 1,2\,m × 1,5\,m rematado por arriba con un semicírculo. ¿Cuánto cuesta el cristal si el m$^2$ vale 45\,€?\par\smallskip\begin{center}\begin{tikzpicture}[x=1.6cm,y=1.6cm,baseline=(current bounding box.center),line join=round,line cap=round]
\draw[fill=teal!12,draw=teal,thick] (0.000,0.000) -- (1.200,0.000) -- (1.200,1.500) -- (0.000,1.500) -- cycle;
\draw[fill=teal!12,draw=teal,thick] (0.600,1.500) -- (1.200,1.500) arc[start angle=0,end angle=180,radius=0.6] -- cycle;
\draw[dashed,muted] (0.000,1.500) -- (1.200,1.500);
\node[font=\small,below] at (0.600,0.000) {1,2 m};
\node[font=\small,right] at (1.200,0.750) {1,5 m};
\end{tikzpicture}\end{center}


% =====================================================
\seccio{10}{El teorema de Pitágoras}

\ex{1} En cada triángulo rectángulo, di qué lados son los catetos y cuál es la hipotenusa, y escribe la relación que da el teorema de Pitágoras:\par\smallskip\begin{center}\begin{tikzpicture}[x=0.45cm,y=0.45cm,baseline=(current bounding box.center),line join=round,line cap=round]
\draw[fill=accent!22,draw=accent,thick] (0.000,0.000) -- (4.000,0.000) -- (0.000,3.000) -- cycle;
\draw[thin] (0.489,0.000) -- (0.489,0.489) -- (0.000,0.489);
\node[font=\small,below] at (2.000,0.000) {$x$};
\node[font=\small,left] at (0.000,1.500) {$y$};
\node[font=\small,above right] at (2.000,1.500) {$z$};
\end{tikzpicture}\qquad\qquad \begin{tikzpicture}[x=0.45cm,y=0.45cm,baseline=(current bounding box.center),line join=round,line cap=round]
\draw[fill=accent!22,draw=accent,thick] (0.000,0.000) -- (5.000,0.000) -- (1.800,2.400) -- cycle;
\draw[thin] (1.507,2.009) -- (1.898,1.716) -- (2.191,2.107);
\node[font=\small,below] at (2.500,0.000) {$m$};
\node[font=\small,above left] at (0.900,1.200) {$n$};
\node[font=\small,above right] at (3.400,1.200) {$p$};
\end{tikzpicture}\qquad\qquad \begin{tikzpicture}[x=0.45cm,y=0.45cm,baseline=(current bounding box.center),line join=round,line cap=round]
\draw[fill=accent!22,draw=accent,thick] (0.000,0.000) -- (0.000,4.000) -- (3.000,4.000) -- cycle;
\draw[thin] (0.000,3.511) -- (0.489,3.511) -- (0.489,4.000);
\node[font=\small,left] at (0.000,2.000) {$r$};
\node[font=\small,above] at (1.500,4.000) {$s$};
\node[font=\small,below right] at (1.500,2.000) {$t$};
\end{tikzpicture}\end{center}\recorda{En un triángulo rectángulo, los \textbf{catetos} forman el ángulo recto y la \textbf{hipotenusa} es el lado opuesto, el más largo. \textbf{Teorema de Pitágoras:} $c^2=a^2+b^2$.}

\ex{1} Calcula la hipotenusa de un triángulo rectángulo de catetos:\begin{tasks}(4)
  \task 3\,cm y 4\,cm
  \task 6\,m y 8\,m
  \task 5\,cm y 12\,cm
  \task 9\,dm y 12\,dm
\end{tasks}\recorda{Hipotenusa: $c=\sqrt{a^2+b^2}$ (sumamos los cuadrados y hacemos la raíz). \ Cateto: $b=\sqrt{c^2-a^2}$ (restamos y hacemos la raíz).}

\ex{1} Calcula el cateto que falta en cada triángulo rectángulo:\begin{tasks}(2)
  \task hipotenusa 10\,cm, cateto 6\,cm
  \task hipotenusa 13\,m, cateto 12\,m
  \task hipotenusa 17\,cm, cateto 15\,cm
  \task hipotenusa 25\,cm, cateto 7\,cm
\end{tasks}

\ex{2} Calcula el lado que falta y redondea el resultado a las centésimas:\begin{tasks}(2)
  \task catetos 2\,cm y 5\,cm
  \task catetos 4\,cm y 6\,cm
  \task hipotenusa 8\,cm, cateto 3\,cm
  \task hipotenusa 10\,cm, cateto 7\,cm
\end{tasks}

\ex{2} Averigua si son rectángulos los triángulos que tienen estos lados:\begin{tasks}(4)
  \task 6, 8 y 10\,cm
  \task 5, 7 y 9\,cm
  \task 8, 15 y 17\,cm
  \task 10, 11 y 15\,cm
\end{tasks}\recorda{Compara el cuadrado del lado más largo con la suma de los cuadrados de los otros dos. Si son iguales, el triángulo es rectángulo.}

\ex{2} Una escalera de 5\,m está apoyada en una pared, con el pie a 1,4\,m de la pared. ¿A qué altura llega?

\ex{2} Un barco sale del puerto y navega 12\,km hacia el norte y después 9\,km hacia el este. ¿A qué distancia en línea recta está del puerto?

\ex{3} En la cuadrícula, cada cuadradito mide 1\,cm de lado. Calcula la distancia entre los puntos $A$ y $B$.\par\smallskip\begin{center}\begin{tikzpicture}[x=0.3cm,y=0.3cm,baseline=(current bounding box.center),line join=round,line cap=round]
\draw[muted!80,line width=0.5pt,step=1] (0.000,0.000) grid (8.000,10.000);
\draw[fill=accent,fill opacity=0.22,draw=accent,thick] (1.000,1.000) -- (7.000,1.000) -- (7.000,9.000) -- cycle;
\draw[thin] (6.267,1.000) -- (6.267,1.733) -- (7.000,1.733);
\fill (1.000,1.000) circle (1.5pt);
\fill (7.000,9.000) circle (1.5pt);
\node[font=\small,below left] at (1.000,1.000) {$A$};
\node[font=\small,above right] at (7.000,9.000) {$B$};
\end{tikzpicture}\end{center}

\ex{3} Para sujetar un poste de 12\,m de alto se atan cuatro cables desde la punta hasta el suelo, cada uno a 5\,m del pie del poste. ¿Cuántos metros de cable hacen falta?\par\smallskip\begin{center}\begin{tikzpicture}[x=0.17cm,y=0.17cm,baseline=(current bounding box.center),line join=round,line cap=round]
\draw[thick] (-1.000,0.000) -- (7.000,0.000);
\draw[accent,very thick] (0.000,0.000) -- (0.000,12.000);
\draw[thick] (0.000,12.000) -- (5.000,0.000);
\draw[thin] (1.294,0.000) -- (1.294,1.294) -- (0.000,1.294);
\node[font=\small,left] at (0.000,6.000) {12 m};
\node[font=\small,below] at (2.500,0.000) {5 m};
\node[font=\small,right] at (2.500,6.000) {cable};
\end{tikzpicture}\end{center}


% =====================================================
\seccio{11}{Pitágoras para calcular áreas y perímetros}

\ex{1} Calcula la diagonal de:\begin{tasks}(3)
  \task un rectángulo de 8\,cm × 6\,cm
  \task un rectángulo de 24\,m × 7\,m
  \task un cuadrado de 4\,cm de lado
\end{tasks}\recorda{Busca dentro de la figura un \textbf{triángulo rectángulo} en el que la medida que te falta sea un lado y del que conozcas los otros dos. Aplica Pitágoras y después calcula el área o el perímetro.}

\ex{1} Calcula la altura, el área y el perímetro de este triángulo isósceles:\par\smallskip\begin{center}\begin{tikzpicture}[x=0.2cm,y=0.2cm,baseline=(current bounding box.center),line join=round,line cap=round]
\draw[fill=teal!12,draw=teal,thick] (-8.000,0.000) -- (8.000,0.000) -- (0.000,6.000) -- cycle;
\draw[dashed,muted] (0.000,0.000) -- (0.000,6.000);
\draw[thin] (1.100,0.000) -- (1.100,1.100) -- (0.000,1.100);
\node[font=\small,below] at (0.000,0.000) {16 cm};
\node[font=\small,above left] at (-4.000,3.000) {10 cm};
\node[font=\small,above right] at (4.000,3.000) {10 cm};
\node[font=\small,right] at (0.000,3.000) {$h$};
\end{tikzpicture}\end{center}

\ex{2} Las diagonales de un rombo miden 18\,cm y 24\,cm. Calcula su lado, su perímetro y su área.

\ex{2} Calcula el perímetro y el área de este trapecio rectángulo:\par\smallskip\begin{center}\begin{tikzpicture}[x=0.17cm,y=0.17cm,baseline=(current bounding box.center),line join=round,line cap=round]
\draw[fill=teal!12,draw=teal,thick] (0.000,0.000) -- (12.000,0.000) -- (7.000,12.000) -- (0.000,12.000) -- cycle;
\draw[dashed,muted] (7.000,0.000) -- (7.000,12.000);
\draw[thin] (1.294,0.000) -- (1.294,1.294) -- (0.000,1.294);
\draw[thin] (8.294,0.000) -- (8.294,1.294) -- (7.000,1.294);
\node[font=\small,below] at (6.000,0.000) {12 cm};
\node[font=\small,above] at (3.500,12.000) {7 cm};
\node[font=\small,left] at (0.000,6.000) {12 cm};
\node[font=\small,right] at (9.500,6.000) {$c$};
\end{tikzpicture}\end{center}

\ex{2} Un trapecio isósceles tiene las bases de 16\,cm y 4\,cm y los lados iguales de 10\,cm. Calcula su altura, su área y su perímetro.

\ex{2} Calcula la apotema y el área de un hexágono regular de 4\,cm de lado.\par\smallskip\begin{center}\begin{tikzpicture}[x=0.55cm,y=0.55cm,baseline=(current bounding box.center),line join=round,line cap=round]
\draw[fill=teal!12,draw=teal,thick] (-1.000,-1.732) -- (1.000,-1.732) -- (2.000,0.000) -- (1.000,1.732) -- (-1.000,1.732) -- (-2.000,0.000) -- cycle;
\draw[dashed,muted] (0.000,0.000) -- (0.000,-1.732);
\draw[dashed,muted] (0.000,0.000) -- (1.000,-1.732);
\draw[thin] (0.000,-1.332) -- (0.400,-1.332) -- (0.400,-1.732);
\fill (0.000,0.000) circle (1.5pt);
\node[font=\small,below] at (0.000,-1.732) {4 cm};
\node[font=\small,left] at (0.000,-0.850) {$a$};
\end{tikzpicture}\end{center}

\ex{3} Un triángulo equilátero tiene 18\,cm de perímetro. Calcula su altura y su área.

\ex{3} Un cuadrado tiene los cuatro vértices sobre una circunferencia de 4\,cm de radio. Calcula el área de la zona coloreada ($\pi\approx3,14$).\par\smallskip\begin{center}\begin{tikzpicture}[x=0.32cm,y=0.32cm,baseline=(current bounding box.center),line join=round,line cap=round]
\draw[fill=accent!22,draw=accent,thick] (0.000,0.000) circle (4);
\draw[fill=white,draw=teal,thick] (4.000,0.000) -- (0.000,4.000) -- (-4.000,0.000) -- (0.000,-4.000) -- cycle;
\draw[dashed,muted] (-4.000,0.000) -- (4.000,0.000);
\fill (0.000,0.000) circle (1.5pt);
\node[font=\small,above] at (2.000,0.000) {4 cm};
\end{tikzpicture}\end{center}

\ex{3} Calcula el perímetro y el área de la fachada de esta casita, con un tejado en forma de triángulo isósceles:\par\smallskip\begin{center}\begin{tikzpicture}[x=0.3cm,y=0.3cm,baseline=(current bounding box.center),line join=round,line cap=round]
\draw[fill=teal!12,draw=teal,thick] (0.000,0.000) -- (8.000,0.000) -- (8.000,5.000) -- (0.000,5.000) -- cycle;
\draw[fill=accent!22,draw=accent,thick] (0.000,5.000) -- (8.000,5.000) -- (4.000,8.000) -- cycle;
\draw[dashed,muted] (4.000,5.000) -- (4.000,8.000);
\draw[thin] (4.733,5.000) -- (4.733,5.733) -- (4.000,5.733);
\node[font=\small,below] at (4.000,0.000) {8 m};
\node[font=\small,right] at (8.000,2.500) {5 m};
\node[font=\small,left] at (4.000,6.500) {3 m};
\end{tikzpicture}\end{center}

\vfill
\begin{center}\color{muted}\small\faCheckCircle[regular]\ \ ¡Ánimo! Si te cuesta algún ejercicio, repasa los apuntes interactivos de la unidad.\end{center}

\end{document}
